\Needspace{6\baselineskip}
\subsection{Conclusions}
The prototype combines external panel sensing, an embedded button controller, and a local wireless application around the Rheem EW1-50RP. Mounted calibration produced distinguishable red, green, and unlit RGBC responses, including small nonzero dark readings. A supervised tape-mounted power-button trial registered a press at a configured angle of 29 degrees. These results provide concrete optical and mechanical references for the assembled device.

Separating current temperature from setpoint was essential to the control logic because the panel flashes the selected value during adjustment. The control design likewise distinguishes standby from power-off without treating an OCR failure as an off observation. Recorded software checks exercised these decisions and the circular selection range, while MQTT and image-recognition failures limit the demonstrated end-to-end operation. The OLED mirrors an accepted panel reading rather than sensing temperature independently, and the simulated flow channel does not establish water use at the heater.

\Needspace{6\baselineskip}
\subsection{Recommendations}
Calibration should be repeated after changing the optical cover or mechanical mounting, and the settling interval should be matched to the observed panel behavior. Priority should be given to confirming MQTT delivery, numerical recognition, and button press-and-return behavior in the same supervised sequence. Future physical flow calibration in a separate test circuit could support volume comparisons without intervening in the heater plumbing. Voice interaction remains a possible extension of the application API.
